% ============================================================
% Unified Methodology -- DRAFT
% ============================================================

\documentclass[11pt, a4paper]{article}

\usepackage[margin=2.2cm]{geometry}
\usepackage{booktabs}
\usepackage{parskip}
\usepackage{enumitem}
\usepackage{xcolor}
\usepackage{fancyhdr}
\usepackage{array}

\pagestyle{fancy}
\fancyhf{}
\rhead{\textcolor{gray}{\small DRAFT}}
\lhead{\textcolor{gray}{\small Unified Methodology}}
\cfoot{\thepage}

\title{\textbf{Unified Methodology for Measuring AI Coding Agent Adoption on GitHub}}
\author{Manuel}
\date{October 2026}

\begin{document}
\maketitle
\thispagestyle{fancy}

% ============================================================
% PAGE 1 -- REPOSITORY COLLECTION
% ============================================================

\section{Repository Collection}

\subsection*{Step 1a -- Pulling Repositories from GitHub}

Repositories are retrieved from the GitHub Search API
(\texttt{/search/repositories}).
Because the API returns at most 1,000 results per query, the search window is
partitioned by creation date (\texttt{created:YYYY-MM-DD..YYYY-MM-DD}) and
iterated over the full date range. Windows with more than 1,000 results are
split in half recursively until each window fits within the cap.

\begin{table}[h]
\centering
\caption{Filters agreed across all team stacks (Step 1a).}
\begin{tabular}{@{}lll@{}}
\toprule
\textbf{Filter} & \textbf{Value} & \textbf{Rationale} \\
\midrule
Stars            & $\geq 4$              & Removes placeholder and abandoned repos \\
Commits          & $\geq 50$             & Ensures meaningful development activity \\
Fork             & \texttt{false}        & Original repositories only \\
Creation date    & $\geq$ 2025-02-24     & Repos born in the AI coding agent era \\
Language         & Stack primary language & JS or TS for this study; differs per stack \\
Language bytes   & $\geq 175{,}000$ (175 KB) & Confirms language is primary, not incidental \\
\bottomrule
\end{tabular}
\end{table}

The \textbf{ecosystem start date} is the date when AI coding agent tooling
became broadly accessible to ordinary developers in the target ecosystem, not
when the first tool was ever released. All repositories created before this
date are excluded, since the tooling was not yet available at the time of
their creation.

\subsection*{Step 1b -- Ecosystem Categorization}

After collection, each repository is categorized by its sub-ecosystem type.
This is done via a single GraphQL query per batch (15 repos at a time) that
checks whether specific files exist at the repository root. No file contents
are read. The detected files determine the category.

\begin{table}[h]
\centering
\caption{Ecosystem categorization logic -- general form.}
\begin{tabular}{@{}lll@{}}
\toprule
\textbf{Category} & \textbf{Detected by} & \textbf{Example (JS/TS stack)} \\
\midrule
Sub-ecosystem A   & Config file A present             & \texttt{package-lock.json} (npm) \\
Sub-ecosystem B   & Config file B present             & \texttt{yarn.lock} (Yarn) \\
Sub-ecosystem C   & Config file C present             & \texttt{pnpm-lock.yaml} (pnpm) \\
Sub-ecosystem D   & Config file D present             & \texttt{bun.lockb} (Bun) \\
No lockfile       & Root config present, no lockfile  & \texttt{package.json} only \\
Alternative runtime & Runtime config present          & \texttt{deno.json} (Deno) \\
Vanilla           & No recognized config              & Plain JS/TS files \\
\bottomrule
\end{tabular}
\end{table}

\textbf{JS/TS example.} Seven files are checked per repository:
\texttt{package.json}, \texttt{package-lock.json}, \texttt{yarn.lock},
\texttt{pnpm-lock.yaml}, \texttt{bun.lockb}, \texttt{deno.json}, and
\texttt{deno.jsonc}. Priority when multiple lockfiles exist: npm $>$ yarn
$>$ pnpm $>$ bun.

\textbf{Python example.} The equivalent files to check would be:
\texttt{requirements.txt}, \texttt{Pipfile.lock}, \texttt{poetry.lock},
\texttt{setup.py}, \texttt{pyproject.toml}, and \texttt{conda.yml}. The
category labels would map to pip, Pipenv, Poetry, conda, and so on.

The output is a single corpus file with an \texttt{ecosystem\_type} column
added to every repository record. No repositories are dropped at this step;
all are kept with their assigned category label.

\newpage
% ============================================================
% PAGE 2 -- COUNTRY DETECTION
% ============================================================

\section{Owner Location Inference}

Location on GitHub is attached to user profiles, not to repositories. For
each unique owner login in the corpus, a country is inferred in two passes.

\subsection*{Pass 1 -- Fast Lookup}

The GitHub GraphQL API is queried in batches of 50 owners per request to
fetch the \texttt{location} profile field. The raw string is matched against
a keyword lookup table of country names, major cities, and common
abbreviations in multiple languages (English and Arabic are covered; others
can be added). Matching uses substring search with the longest key checked
first to avoid partial matches.

Owners whose field is empty, unrecognisable, or contains only generic terms
(e.g., ``Earth'', ``Remote'') are passed to Pass 2.

\subsection*{Pass 2 -- Deep Profile Inference}

One GraphQL query per unresolved owner fetches additional profile fields. Six
signals are tried in order; inference stops at the first successful match.

\begin{table}[h]
\centering
\caption{Signal cascade in Pass 2.}
\begin{tabular}{@{}clll@{}}
\toprule
\textbf{Priority} & \textbf{Signal} & \textbf{Field fetched} & \textbf{Confidence} \\
\midrule
1 & Bio keyword match    & \texttt{bio}                        & Medium \\
2 & Company keyword match & \texttt{company}                   & Medium \\
3 & Company domain TLD   & \texttt{company} (domain part)      & Medium \\
4 & Email domain TLD     & \texttt{email}                      & Medium \\
5 & Website TLD          & \texttt{websiteUrl}                 & Low    \\
6 & Organisation location & \texttt{organizations.nodes.location} & Medium \\
7 & Network majority     & \texttt{followers} + \texttt{following} ($\geq 3$ located, $\geq 40\%$ agreement) & Low \\
\bottomrule
\end{tabular}
\end{table}

Country-code top-level domains (e.g., \texttt{.de} for Germany, \texttt{.eg}
for Egypt, \texttt{.br} for Brazil) are used for signals 3, 4, and 5.
Owners that remain unresolved after all seven signals are labelled
\texttt{not\_specified} and are excluded from country-level analysis.

\subsection*{Combined Output}

Pass 1 and Pass 2 results are merged into \texttt{owner\_locations\_all.csv}.
This file maps every owner login to a country name and a confidence level.
It is used by the geographic analysis phase to assign a country to each
repository via its owner.

\newpage
% ============================================================
% PAGE 3 -- AGENTIC TRACES DETECTION
% ============================================================

\section{AI Adoption Detection}

Each repository is classified as \textbf{agentic}, \textbf{AI-assisted}, or
\textbf{none} using a three-pass layered pipeline. A repository is classified
by the first pass that finds a signal and never re-processed by later passes.
This keeps the total API cost low while still applying expensive signals where
needed.

\subsection*{Pass 1 -- Configuration File Check}

Repositories are checked in batches of 10 via GraphQL for the presence of
known AI agent configuration files at the repository root. These files are
hard positives: they are created deliberately by developers configuring an
agent.

\begin{table}[h]
\centering
\caption{AI agent configuration files checked in Pass 1.}
\begin{tabular}{@{}ll@{}}
\toprule
\textbf{File} & \textbf{Agent} \\
\midrule
\texttt{CLAUDE.md}                       & Claude Code \\
\texttt{AGENTS.md}                       & Claude Code / Generic \\
\texttt{.cursorrules}                    & Cursor \\
\texttt{.cursorignore}                   & Cursor \\
\texttt{.aider.conf.yml}                 & Aider \\
\texttt{.roomodes}                       & Roo Code \\
\texttt{.windsurfrules}                  & Windsurf \\
\texttt{GEMINI.md}                       & Gemini CLI \\
\texttt{.github/copilot-instructions.md} & GitHub Copilot \\
\texttt{.continue/config.json}           & Continue \\
\texttt{.codeium/config.json}            & Codeium \\
\texttt{AGENTS.json}                     & OpenAI Codex \\
\bottomrule
\end{tabular}
\end{table}

Any match immediately labels the repository \textbf{agentic}.

\subsection*{Pass 2 -- Commit Trailer and PR Branch Scan}

The 100 most recent commits and 50 most recent pull requests are fetched.
Commit messages are scanned for co-authorship trailer patterns. PR head
branch names are scanned for agent-specific prefixes.

\begin{table}[h]
\centering
\caption{Representative patterns checked in Pass 2.}
\begin{tabular}{@{}ll@{}}
\toprule
\textbf{Pattern} & \textbf{Agent} \\
\midrule
\texttt{Co-authored-by: ...claude...anthropic} & Claude Code \\
\texttt{Co-authored-by: ...github.copilot}     & GitHub Copilot \\
\texttt{Co-authored-by: ...devin}              & Devin \\
\texttt{Co-authored-by: ...cursor}             & Cursor \\
\texttt{Co-authored-by: ...aider}              & Aider \\
\texttt{Co-authored-by: ...windsurf}           & Windsurf \\
\texttt{Co-authored-by: ...gemini}             & Gemini \\
\texttt{noreply@anthropic.com} (author email)  & Claude Code \\
PR branch prefix \texttt{copilot/}            & GitHub Copilot \\
PR branch prefix \texttt{codex/}              & OpenAI Codex \\
PR branch prefix \texttt{devin/}              & Devin \\
\bottomrule
\end{tabular}
\end{table}

\subsection*{Pass 3 -- Deep Scan}

Repositories not classified by Pass 1 or Pass 2 receive a comprehensive scan:
up to 300 commits (message and author email), up to 100 pull requests (body
text, labels, and branch name), and the package dependency file
(\texttt{package.json} for npm, or the equivalent for other ecosystems).

\subsection*{Classification Rule (all passes)}

\begin{table}[h]
\centering
\caption{Classification thresholds applied after each pass.}
\begin{tabular}{@{}lll@{}}
\toprule
\textbf{Condition} & \textbf{Label} \\
\midrule
$\geq 3$ AI signals, OR signals from $\geq 2$ distinct agents & \textbf{agentic} \\
1 or 2 AI signals from a single agent                         & \textbf{AI-assisted} \\
No signal found in any pass                                   & \textbf{none} \\
\bottomrule
\end{tabular}
\end{table}

After all three passes complete, results are merged in priority order (Pass 1
supersedes Pass 2, which supersedes Pass 3) and joined with the corpus to
produce \texttt{corpus\_detected\_full.csv}.

\end{document}
